\documentclass[10pt,a4paper]{amsart}
\usepackage[margin=19mm]{geometry}
\usepackage{amsmath,amssymb,mathtools}
\usepackage[scr=rsfs,cal=euler]{mathalfa}
\usepackage{xfrac}
\usepackage[unicode,hidelinks,bookmarks=false]{hyperref}
\newcommand{\ie}{i.e.\ }
\newcommand{\cf}{cf.\ }
\newcommand{\resp}{respectively\ }
\newcommand{\Subsection}{\subsection}
\providecommand{\nobreakdash}{\nobreak\hbox{-}}
\setlength{\emergencystretch}{2em}
\multlinegap0pt
\usepackage{amssymb}
\usepackage{enumerate}
\usepackage{xcolor}
\usepackage{mathtools}
\usepackage{tikz}
\newtheorem{lemma}{Lemma}
\newcommand{\R}{\mathbb{R}}
\newcommand{\cE}{\mathcal{E}}
\newcommand{\dd}{\mathrm{d}}
\newcommand{\red}{}
\title{From simplex slicing to sharp reverse H\"older inequalities}
\author{James Melbourne, Michael Roysdon, Colin Tang, Tomasz Tkocz}
\date{Source: arXiv:2505.00944v2 (CC BY 4.0)}
\begin{document}
\maketitle

\noindent\emph{Source and licence.} From simplex slicing to sharp reverse H\"older inequalities. Version arXiv:2505.00944v2 (CC BY 4.0), distributed by arXiv under the Creative Commons Attribution 4.0 International licence, http://creativecommons.org/licenses/by/4.0/, as recorded in the arXiv licence metadata for this exact version and shown on the abstract page https://arxiv.org/abs/2505.00944v2. Full licence notice: \texttt{licenses/arxiv-2505.00944-CC-BY-4.0.txt}.\par\smallskip

\noindent\textit{Editorial scope.} Counted unit: the article's lemma lm:3crossings (source0.tex lines 527--529). The article's complete original proof of it is delivered with it (source0.tex lines 530--582, one \texttt{proof} environment of 53 lines). No mathematics is written, repaired or re-derived: the statement and the proof are the article's own lines, in the article's own notation, and no retained line is substituted or re-flowed.  Retained uncounted as the source-local context the proof needs: lines 329--335; lines 358--360; lines 362--364.  Nothing was omitted from any retained span, so the package carries no omission record.  No retained line contains a citation, so the wrapper carries no bibliography and no bibliography entry.  No retained line contains a figure environment, an image-inclusion command or drawing code, so no image is delivered and the package declares no asset dependency.  Editorial formatting only, in addition: an AMS \texttt{amsart} preamble that restates the article's own declarations and macros used by the retained lines, namely the environments \texttt{lemma} on separate counters, together with the standard packages those lines need.  The retained environments print in this document as Lemma 1 in order of appearance, whereas the article prints the same items with the numbering of its own full document.  The complete native source of the article as distributed in the arXiv e-print for this exact version is bundled unchanged as \texttt{sources/177461/source0.tex}.
\begin{lemma}\label{lm:simple-crossings}
Let $h \colon (0, +\infty) \to \R$ be an integrable function. Suppose $\int_0^{+\infty} h(x) \dd x = 0$. If  $h$ has exactly one weak sign change with sign pattern $(0-,0+)$, and if $h$ is nonzero on a set of positive measure, then
\[ 
\int_0^{+\infty} h(x)x \dd x > 0
 \]
(provided the integrability of $h(x)x$). 
\end{lemma}

\begin{equation}\label{eq:def-Xab}
X_{a,b} = a(\cE-1) - b(\cE'-1).
 \end{equation}

\begin{equation}\label{eq:gab}
g_{a,b}(x) = \begin{cases}\frac{1}{a+b}e^{-(x+a-b)/a}, & x \geq -(a-b), \\ \frac{1}{a+b}e^{(x+a-b)/b}, & x < -(a-b) \end{cases}
\end{equation}

\begin{lemma}\label{lm:3crossings}
Let $f_t\colon[0,{+\infty}) \to [0,{+\infty})$ be the density of $|\bar\cE_t|$. For every $0<t<1$, each of the functions $f_t - f_0$ and $f_1 - f_t$ changes sign exactly $3$ times on $(0, {+\infty})$ and has sign pattern $(+,-,+,-)$.
\end{lemma}

\begin{proof}
In the notation of \eqref{eq:def-Xab}, we have $\cE_t = X_{1,t}$, so recalling \eqref{eq:gab}, $\cE_t$ has density $g_{1,t}$, thus $|\cE_t|$ has density on $(0,{+\infty})$ at $x > 0$ equal to
\begin{align*}
\rho_t(x) = g_{1,t}(x) + g_{1,t}(-x) &= \frac{1}{1+t}e^{-(x+1-t)} +  \frac{1}{1+t}\begin{cases}
e^{x-(1-t)}, & x \leq 1-t, \\ e^{-x/t+(1-t)/t}, & x > 1-t,
\end{cases}\\
&=\frac{e^{t-1}}{1+t}\begin{cases}
e^{x}+e^{-x}, & x \leq 1-t, \\ e^{-x} + e^{-x/t + 1/t-t}, & x > 1-t.
 \end{cases}
 \end{align*}
When $t = 0$ (following the convention in \eqref{eq:gab} that $e^{-x/t+(1-t)/t}$ is replaced by $0$ identically on $(1,{+\infty})$),
\[ 
\rho_0(x) = e^{-1}\begin{cases}
e^{x}+e^{-x}, & x \leq 1, \\ e^{-x}, & x > 1.
 \end{cases}
 \]
Plainly 
\[ 
\rho_1(x) = e^{-x}.
 \]
 Note that for every $0 < t \leq 1$, $\rho_t$ is continuous, whereas $\rho_0$ has a jump at $x = 1$ and is continuous elsewhere.

The density $f_t$ of $|\bar\cE_t| = \mu_t^{-1}|\cE_t|$ is then $f_t(x) = \mu_t\rho_t(\mu_tx)$. 

Fix $0 < t < 1$. First let us handle the function 
\[
h_1(x) = f_1(x) - f_t(x) = \rho_1(x) - \mu_t\rho_t(\mu_tx).
\]
Since $\int_0^{+\infty} h_1 = 0$, and $h_1$ is not identically $0$, it must change sign at least once; additionally $\int_0^{+\infty} h_1(x)x \dd x = 0$, so $h_1$ need change sign at least twice on $(0, {+\infty})$, by Lemma \ref{lm:simple-crossings}.
It is evident that the ratio $r_1(x) = f_t(x)/f_1(x) = \mu_t\rho_t(\mu_t x)e^x$ is convex on each of the intervals $[0, \alpha_t]$ and $[\alpha_t, {+\infty})$, where $\alpha_t = \mu_t^{-1}(1-t)$ denotes the point where the formula defining $f_t(x)$ changes. Moreover, $r_1(x)$ is strictly increasing on $[0, \alpha_t]$ (as the product of two strictly increasing functions) with $r_1(0) = \mu_t\rho_t(0) = \mu_t^2 < 1$. It is also clear that $r_1(x) \to {+\infty}$ as $x \to {+\infty}$. If $r_1(\alpha_t) \leq 1$, then $r_1(x)=1$ would have no solutions on $(0,\alpha_t)$ and exactly $1$ solution on $(\alpha_t,{+\infty})$, resulting in $f_1 - f_t$ having at most one sign change, ruled out earlier. Therefore, $r_1(\alpha_t) > 1$. Consequently, $r_1(x)=1$ has exactly one solution on $(0,\alpha_t)$ (by monotonicity) and exactly $2$ solutions on $(\alpha_t,{+\infty})$ (by convexity), resulting in $f_1 - f_t$ having exactly three sign changes and the sign pattern $(+,-,+,-)$ (because $r_1(0) < 1$).

Now we handle the function
\[
h_0(x) = f_t(x) - f_0(x) = \mu_t\rho_t(\mu_tx) - \mu_0\rho_0(\mu_0x).
\]
As before, it necessarily has at least two sign changes.

We will frequently use that $\mu_t > \mu_0$. Recall $\rho_0(\mu_0x)$ has a jump discontinuity at $x_0 = \mu_0^{-1}$ and the formula for $\rho_t(\mu_tx)$ changes at $x = \alpha_t = \mu_t^{-1}(1-t) < \mu_0^{-1} = x_0$. It will be convenient to split the analysis into the consecutive intervals: $[0, \alpha_t]$, $(\alpha_t, x_0]$ and $(x_0,+ \infty)$.

On $[0, \alpha_t]$, we have $\frac12 h_0(x) = \mu_t\frac{e^{t-1}}{1+t}\cosh(\mu_tx) - \mu_0e^{-1}\cosh(\mu_0x)$. Since $\mu_t > \mu_0$ and $\frac{e^{t-1}}{1+t} > e^{-1}$, we have $h_0 > 0$ on $[0, \alpha_t]$.

On $[\alpha_t, x_0]$, $f_t$ is now strictly decreasing, whereas $f_0(x) = \mu_0^2\cosh(\mu_0 x)$ is still strictly increasing, so $h_0$ changes sign at most once on this interval.

On $(x_0, {+\infty})$, the ratio 
\[
r_0(x) = \frac{f_t(x)}{f_0(x)} = \frac{e\mu_t}{\mu_0}\rho_t(\mu_tx)e^{\mu_0x} = \frac{\mu_t}{\mu_0}\frac{e^t}{1+t}\left(e^{-(\mu_t-\mu_0)x} +e^{-(\mu_t/t-\mu_0)x + 1/t-t} \right)
\]
is clearly strictly decreasing with $r_0(x) \to 0$ as $x \to {+\infty}$. In particular, $h_0$ can have at most one sign change here, and is eventually negative. 

Function $h_0$ admits one more sign change, at the discontinuity point $x_0$, so at most $3$ sign changes on $(0, {+\infty})$, and minimum $2$ sign changes. As remarked, near $0$, $h_0$ is positive, and eventually negative, thus it {\red needs to} change sign an odd number of times. Thus we get exactly $3$ sign changes. Since $f_t(x) < f_0(x)$ as $x \to {+\infty}$, the sign pattern is $(+,-,+,-)$.

%Moreover, $r_0(x_0-) < r_0(x_0+)$ because $f_0$ jumps down at $x_0$. If $r_0(x_0+) \leq 1$, then there is no sing change on $[x_t, {+\infty})$, so $h_0$ would have only one sign change, a contradiction. Consequently, $r_0(x_0+) > 1$ and $h_0$ changes sign exactly once on $(x_0, {+\infty})$. If $r_0(x_0+) \geq 1$, i.e. $f_t(x_0) \geq f_0(x_0)$ resulting with $f_t > f_0$ on the whole interval $(0, x_0)$ (as established above, $f_t > f_0$ on $[0,y_t]$ and then on $[y_t, x_0]$, $f_t$ decreases to $f_t(x_0)$ staying above $f_0$ which increases). This would give again only one sign change overall. Thus $r_0(x_0+) < 1$, giving a sign change at $x_0$ as well as, one sign change on $(0, x_0)$. Thus we get exactly $3$ sign changes. Since $f_t(x) < f_0(x)$ as $x \to {+\infty}$, the sign pattern is $(+,-,+,-)$.
\end{proof}

\end{document}
